\documentclass[10pt,a4paper]{amsart}
\usepackage[margin=19mm]{geometry}
\usepackage{amsmath,amssymb,mathtools}
\usepackage[scr=rsfs,cal=euler]{mathalfa}
\usepackage{xfrac}
\usepackage[unicode,hidelinks,bookmarks=false]{hyperref}
\newcommand{\ie}{i.e.\ }
\newcommand{\cf}{cf.\ }
\newcommand{\resp}{respectively\ }
\newcommand{\Subsection}{\subsection}
\providecommand{\nobreakdash}{\nobreak\hbox{-}}
\setlength{\emergencystretch}{2em}
\multlinegap0pt
\usepackage{bm}
\usepackage{bbm}
\usepackage{dsfont}
\usepackage{xspace}
\usepackage{cancel}
\usepackage{mathtools}
\usepackage{amsthm}
\makeatletter
\@ifundefined{theorem}{\newtheorem{theorem}{Theorem}}{}
\@ifundefined{proposition}{\newtheorem{proposition}[theorem]{Proposition}}{}
\makeatother
\makeatletter
\expandafter\ifx\csname enum\endcsname\relax
\newenvironment{enum}{\begin{enumerate}
		\renewcommand{\labelenumi}{(\roman{enumi})}}{\end{enumerate}}
\fi
\newcommand{\rr}{{\mathbb{R}}}
\makeatother
\title[Variational properties of the total inverse mean curvature i]{Variational properties of the total inverse mean curvature in the plane under boundary constraints}
\author{Juli\'an Pozuelo, Simone Verzellesi, Giacomo Vianello}
\date{Source: arXiv:2510.25576v1 (CC BY 4.0)}
\begin{document}
\maketitle

\noindent\emph{Source and licence.} Variational properties of the total inverse mean curvature in the plane under boundary constraints. Version arXiv:2510.25576v1 (CC BY 4.0), distributed under the Creative Commons Attribution 4.0 International licence, as recorded in the arXiv licence metadata for this exact version and shown on the abstract page https://arxiv.org/abs/2510.25576v1. Full rights notice: licenses/arxiv-2510.25576-CC-BY-4.0.txt.\par\smallskip

\noindent\textit{Editorial scope.} Counted unit: the Proposition whose source label is \texttt{soluzgenrralestab} in the article (source0.tex lines 1823--1866, source label \texttt{soluzgenrralestab}), delivered together with the article's own complete original proof of it (source0.tex lines 1867--2067, one \texttt{proof} environment of 201 lines). No mathematics is written, repaired or re-derived: statement and proof are the article's own text line for line. Required context retained uncounted: equazionepuntocriticoinlemma (lines 928--930); onlyequation (lines 1820--1822). Every definition, display and label of those lines is retained unchanged. Omitted from this package and not redistributed: 1--927, 931--1819, 2068--2497. Nothing inside a retained excerpt is omitted or altered: every retained body is delivered exactly as the article's lines are written, so no omission receipt of any kind accompanies this package. Citations: the retained excerpts carry no citation command at all, so no bibliography entry of the article is transcribed and no reference is redistributed. Reference adaptation: none was needed and none was made; every reference inside the delivered unit resolves inside it, so no unresolved reference or citation is produced. Printed slips kept literal: no printed word, symbol, hypothesis or number of the counted statement or of its proof was altered. Layout and class: the article's own class is \texttt{amsart} with packages including xcolor, amsfonts, amssymb, amsmath, amscd, amstext, pgfplots, hyperref, inputenc, graphicx, changes, comment, mathrsfs, listings; the wrapper is the lane's pinned \texttt{amsart} editorial wrapper at 10pt on A4, which restates the theorem-like environments the retained text needs and adds the article's own macro definitions that the retained text needs (\texttt{\textbackslash rr}), with extra wrapper packages (bm, bbm, dsfont, xspace, cancel, mathtools). The delivered numbering is the unit's own. Assets: no retained span contains a figure environment, a graphics-inclusion command, a PDF-page inclusion, a table or a TikZ picture; no image file is copied into this package, the package declares no asset dependency and no image file is redistributed with it. Licence: version arXiv:2510.25576v1 of this article is distributed under the Creative Commons Attribution 4.0 International licence, as recorded in the arXiv licence metadata for this exact version and shown on the abstract page https://arxiv.org/abs/2510.25576v1. A full rights notice is delivered at licenses/arxiv-2510.25576-CC-BY-4.0.txt. Characters: every retained excerpt is pure ASCII, so the delivered engine, XeLaTeX with its default Latin Modern text font, reports no missing character.
		there exist a unique $\lambda>0$ and a unique vertically symmetric admissible curve $\gamma=(x,y)$ parametrized by arc-length, such that $\dot y(0)=\dot y(2L)=0$ and        \begin{equation}\label{equazionepuntocriticoinlemma}
			2+\frac{d^2(H^{-2})}{ds^2}=\lambda\qquad\text{ on $[0,2L]$,}
		\end{equation} 

\begin{equation}\label{onlyequation}
\left(r \ddot u\right)''+\mu \left(q \dot u\right)'+pu=0\qquad\text{on }(0,2L),
\end{equation}

\begin{proposition}\label{soluzgenrralestab}
Let $\mu_0=-\frac{x_0}{L+x_0}+2\sqrt{\frac{x_0}{L+x_0}}$. Then, the following facts hold.
\begin{enumerate}
	\item [(i)] Assume that $\mu<\mu_0$.
	%              \begin{equation}\label{mucomplesso}
		%        \deleted{ 0<\mu<-\frac{x_0}{L+x_0}+2\sqrt{\frac{x_0}{L+x_0}}.}
		%   \end{equation}
	Then any solution to \eqref{onlyequation} is a linear combination of the functions    \begin{equation}\label{soluzionicasocomplesso}
		u_1=\cos \alpha\theta\sinh\beta\theta,\qquad
		u_2=\sin \alpha\theta\cosh\beta\theta,\qquad
		u_3= \sin \alpha\theta\sinh\beta\theta,\qquad
		u_4=\cos \alpha\theta\cosh\beta\theta,
	\end{equation}
	where
	\begin{equation}\label{alfabetacomplesso}
		\alpha=\frac{1}{2}\sqrt{\mu+\frac{x_0}{L+x_0}+2\sqrt{\frac{x_0}{L+x_0}}},\qquad\beta=\frac{1}{2}\sqrt{-\mu-\frac{x_0}{L+x_0}+2\sqrt{\frac{x_0}{L+x_0}}}.
	\end{equation}
	\item [(ii)] Assume that $\mu=\mu_0$.
	%     \begin{equation}\label{mucritico}
		%  \deleted{    \mu=-\frac{x_0}{L+x_0}+2\sqrt{\frac{x_0}{L+x_0}}.}
		%   \end{equation}
	Then any solution to \eqref{onlyequation} is a linear combination of the functions     \begin{equation}\label{soluzionicasocritico}
		u_1=\sin\alpha \theta,\qquad
		u_2=\cos \alpha\theta,\qquad
		u_3= \theta\sin \alpha\theta,\qquad
		u_4=\theta\cos \alpha\theta.
	\end{equation}
	\item [(iii)] Assume that $\mu>\mu_0$.
	%\begin{equation}\label{mureale}
	%\deleted{          \mu >-\frac{x_0}%{L+x_0}+2\sqrt{\frac{x_0}{L+x_0}}.}
	%   \end{equation}
Then any solution to \eqref{onlyequation} is a linear combination of the functions
\begin{equation}\label{soluzionicasoreale}
	u_1=\sin\left(\alpha-\gamma\right) \theta,\qquad
	u_2=\sin \left(\alpha+\gamma\right)\theta,\qquad
	u_3= \cos (\alpha-\gamma)\theta,\qquad
	u_4=\cos (\alpha+\gamma)\theta,
\end{equation}
where
\begin{equation}\label{defigamma}
	\gamma=\frac{1}{2}\sqrt{\mu+\frac{x_0}{L+x_0}-2\sqrt{\frac{x_0}{L+x_0}}}.
\end{equation}
\end{enumerate}
\end{proposition}

\begin{proof}
%  \deleted{  First we show that, when }
%         \begin{equation*}
%\deleted{          \mu\neq\frac{x_0}{L+x_0}+2\sqrt{\frac{x_0}{L+x_0}},}
%    \end{equation*}
% \deleted{       $u_1,\ldots,u_4$ are solutions to \eqref{onlyequation}.} 
Fix $z\in\mathbb C$, and set $u:[0,2L]\to\mathbb C$ by $u(s)=\sin z\theta(s)$. Then
\begin{equation*}
\begin{split}
	\left(r  u''\right)''+ \mu\left(q  u'\right)' + pu&= \left(2 H^{-3} (\sin z\theta)''\right)''+ \mu\left(2H^{-1} (\sin z\theta)'\right)' + (2-\lambda)H\sin z\theta\\
	&=\left(2z H^{-3} (H\cos z \theta)'\right)''+ \mu\left(2 z \cos z \theta\right)' + (2-\lambda)H\sin z \theta\\
	%         &\deleted{=\left(2  z H^{-3} \dot H\cos z \theta-2  z ^2H^{-1}\sin z \theta\right)''-2\mu z^2H\sin z\theta + (2-\lambda)H\sin z \theta}\\
	&=\left(- z \left( H^{-2}\right)' \cos z \theta-2  z ^2H^{-1}\sin z \theta\right)''+\left(-2\mu z^2+2-\lambda\right)H\sin z\theta\\
	%        &\deleted{=\left(- z \left( H^{-2}\right)'' \cos z\theta+ z ^2\left( H^{-2}\right)'H \sin z \theta-2 z ^2\left( H^{-1}\right)'\sin z \theta-2 z ^3\cos\theta\right)'}\\
	%&\deleted{\quad+\left(-2\mu z^2+2-\lambda\right)H\sin z\theta}\\
	%        &=\left(-z\left( H^{-2}\right)'' \cos z\theta+2z^2\left( H^{-1}\right)' \sin z\theta-2z^2\left( H^{-1}\right)'\sin z\theta-2z^3\cos\theta\right)'\\
	%     &\quad+\left(-2\mu z^2+2-\lambda\right)H\sin z\theta\\
	&=\left(-z\left( H^{-2}\right)'' \cos z\theta-2z^3\cos\theta\right)'+\left(-2\mu z^2+2-\lambda\right)H\sin z\theta\\
	&\overset{\eqref{equazionepuntocriticoinlemma}}{=} \left(\left(z(2-\lambda)-2z^3\right)\cos z\theta\right)'+\left(-2\mu z^2+2-\lambda\right)H\sin z\theta\\
	&=\left(2z^4-(2\mu+2-\lambda)z^2+(2-\lambda)\right)H\sin z\theta,\\
\end{split}
\end{equation*}
where we used the fact that $\dot\theta=H$ and $H (H^{-2})'=2(H^{-1})'$. In the same way, $v:[0,2L]\to\mathbb C$ defined by $v(s)=\cos z\theta (s)$ satisfies
\begin{equation*}
\left(r  v''\right)''+ \mu\left(q  v'\right)' + pv=\left(2z^4-(2\mu+2-\lambda)z^2+(2-\lambda)\right)H\cos z\theta.
\end{equation*}
In particular, $u$ and $v$ solve \eqref{onlyequation} provided that 
\begin{equation}\label{complexpolinomio}
2z^4-(2\mu+2-\lambda)z^2+(2-\lambda)=0.
\end{equation}
In this case, also the real and imaginary part $u$ and $v$ solve \eqref{onlyequation}. When $\mu<\mu_0$, the four solutions to \eqref{complexpolinomio} are $\pm\alpha\pm\mathrm i \beta$ and $\pm\alpha\mp \mathrm i \beta.$ Hence, the functions $u_1,\ldots,u_4$ given in \eqref{soluzionicasocomplesso} are solutions to \eqref{onlyequation}. When instead $\mu>\mu_0$, the four solutions to \eqref{complexpolinomio} are $\pm\alpha\pm\gamma$ and $\pm\alpha\mp\gamma$, and the functions given in \eqref{soluzionicasoreale} are solutions to \eqref{onlyequation}.
Finally, when $\mu=\mu_0$, \eqref{complexpolinomio} admits two real solutions, $\pm\alpha$. In particular, the functions $u_1$ and $u_2$ given in \eqref{soluzionicasocritico} are solutions. We claim that $u_3$ and $u_4$ are solutions too. To this aim, since $\mu=\mu_0$,
\begin{equation}\label{auxcomputcritic}
4\alpha^2=4\sqrt{\frac{x_0}{L+x_0}}\qquad\text{and}\qquad 2-\lambda+2\mu=4\sqrt{\frac{x_0}{L+x_0}}.
\end{equation}
Therefore,
\begin{equation*}      \begin{split}          \left(r \dot u''\right)''+ \mu&\left(q \dot u'\right)' + pu=\left(2H^{-3}(\theta\sin\alpha\theta)''\right)''+2\mu\left(H^{-1}(\theta\sin\alpha\theta)'\right)'+(2-\lambda)H\theta\sin\alpha\theta\\
	&=\left(2H^{-3}(H\sin\alpha\theta+\alpha H\theta\cos\alpha\theta)'\right)''+2\mu\left(\sin\alpha\theta+\alpha\theta\cos\alpha\theta\right)'+(2-\lambda)H\theta\sin\alpha\theta\\
	%    &\deleted{=\left(2H^{-3}\left(\dot H\sin\alpha\theta+2\alpha H^2\cos\alpha\theta+\alpha \dot H\theta\cos\alpha\theta-\alpha^2 H^2\theta\sin\alpha\theta\right)\right)''}\\
	%    &\deleted{\quad+4\mu\alpha H\cos\alpha\theta-2\mu\alpha^2 H\theta\sin\alpha\theta+(2-\lambda)H\theta\sin\alpha\theta}\\
	&=\left(-\left(H^{-2}\right)'\sin\alpha\theta+4\alpha H^{-1}\cos\alpha\theta-\alpha \left(H^{-2}\right)'\theta\cos\alpha\theta-2\alpha^2 H^{-1}\theta\sin\alpha\theta\right)''\\
	&\quad+4\mu\alpha H\cos\alpha\theta+\left(-2\mu\alpha^2 +(2-\lambda)\right)H\theta\sin\alpha\theta\\
	%    &\deleted{=\left((2-\lambda)\sin\alpha\theta-\alpha H\left(H^{-2}\right)'\cos\alpha\theta+4\alpha\left(H^{-1}\right)'\cos\alpha\theta-4\alpha^2\sin\alpha\theta\right)'}\\
	%    &\deleted{\quad+\left((2-\lambda)\alpha\theta\cos\alpha\theta-\alpha H\left(H^{-2}\right)'\cos\alpha\theta+\alpha^2H\left(H^{-2}\right)'\theta\sin\alpha\theta\right)'}\\
	%   &\deleted{\quad+\left(-2\alpha^2\left(H^{-1}\right)'\theta\sin\alpha\theta-2\alpha^2\sin\alpha\theta-2\alpha^3\theta\cos\alpha\theta\right)'}\\
	%     &\deleted{\quad+4\mu\alpha H\cos\alpha\theta+\left(-2\mu\alpha^2 +(2-\lambda)\right)H\theta\sin\alpha\theta}\\
	&=\left(\left(2-\lambda-6\alpha^2\right)\sin\alpha\theta+\alpha\left(2-\lambda-2\alpha^2\right)\theta\cos\alpha\theta\right)'\\
	&\quad+4\mu\alpha H\cos\alpha\theta+\left(-2\mu\alpha^2 +(2-\lambda)\right)H\theta\sin\alpha\theta\\
	%       &\deleted{=\alpha\left(2-\lambda-6\alpha^2\right)H\cos\alpha\theta+\alpha\left(2-\lambda-2\alpha^2\right)H\cos\alpha\theta-\alpha^2\left(2-\lambda-2\alpha^2\right)H\theta\sin\alpha\theta}\\
	%         &\deleted{\quad+4\mu\alpha H\cos\alpha\theta+\left(-2\mu\alpha^2 +(2-\lambda)\right)H\theta\sin\alpha\theta}\\
	&=2\alpha\left(2-\lambda-4\alpha^2+2\mu\right)H\cos\alpha\theta+\left(2\alpha^4-\left(2\mu+2-\lambda\right)\alpha^2+(2-\lambda)\right)H\theta\sin\alpha\theta\\
	&\overset{\eqref{complexpolinomio}}{=}2\alpha\left(-4\alpha^2+2-\lambda+2\mu\right)H\cos\alpha\theta\\
	&\overset{\eqref{auxcomputcritic}}{=}0,
\end{split}
\end{equation*}
whence $u_3$ is a solution. Similarly,  $u_4$ is a solution too.     
We are left to prove that the above solutions are linearly independent. To this aim, let $A_{ij}=u^{(i-1)}_j(0)$ for $i,j=1,\ldots,4$. Then $u_1,\ldots,u_4$ are linearly independent if and only if $\det A\neq 0$. 
Assume first that $\mu<\mu_0$ and $u_1,\ldots,u_4$ are the functions given by \eqref{soluzionicasocomplesso}.
%Assume first that 
%   \begin{equation*}
%      0<C<-\frac{x_0}{L+x_0}+2\sqrt{\frac{x_0}{L+x_0}}.
%   \end{equation*}
%   Let $x_1,z_2\in\mathbb C$ be as above. Let $\alpha,\beta>0$ be such that $z_1=\alpha-\mathrm i\beta$ and $z_2=\alpha+\mathrm i\beta$. 
We recall that
\begin{equation*}
u_1=\cos \alpha\theta\sinh\beta\theta,\qquad
u_2=\sin \alpha\theta\cosh\beta\theta,\qquad
u_3= \sin \alpha\theta\sinh\beta\theta,\qquad
u_4=\cos \alpha\theta\cosh\beta\theta.
\end{equation*}
Then,
\begin{align*}
\dot u_1&=H\left(-\alpha\sin\alpha\theta\sinh\beta\theta+\beta\cos\alpha\theta\cosh\beta\theta\right),\qquad\dot u_2=H\left(\alpha\cos\alpha\theta\cosh\beta\theta+\beta\sin\alpha\theta\sinh\beta\theta\right),\\
\dot u_3&=H\left(\alpha\cos\alpha\theta\sinh\beta\theta+\beta\sin\alpha\theta\cosh\beta\theta\right),\qquad\dot u_4=H\left(-\alpha\sin\alpha\theta\cosh\beta\theta+\beta\cos\alpha\theta\sinh\beta\theta\right).
\end{align*}
Moreover,
\begin{align*}
\ddot u_1&=\dot H\left(-\alpha\sin\alpha\theta\sinh\beta\theta+\beta\cos\alpha\theta\cosh\beta\theta\right)+H^2\left(-2\alpha\beta\sin\alpha\theta\cosh\beta\theta+\left(\beta^2-\alpha^2\right)\cos\alpha\theta\sinh\beta\theta\right),\\
\ddot u_2&=\dot H\left(\alpha\cos\alpha\theta\cosh\beta\theta+\beta\sin\alpha\theta\sinh\beta\theta\right)+H^2\left(2\alpha\beta\cos\alpha\theta\sinh\beta\theta+\left(\beta^2-\alpha^2\right)\sin\alpha\theta\cosh\beta\theta\right),\\
\ddot u_3&=\dot H\left(\alpha\cos\alpha\theta\sinh\beta\theta+\beta\sin\alpha\theta\cosh\beta\theta\right)+H^2\left(2\alpha\beta\cos\alpha\theta\cosh\beta\theta+\left(\beta^2-\alpha^2\right)\sin\alpha\theta\sinh\beta\theta\right),\\
\ddot u_4&=\dot H\left(-\alpha\sin\alpha\theta\cosh\beta\theta+\beta\cos\alpha\theta\sinh\beta\theta\right)+H^2\left(-2\alpha\beta\sin\alpha\theta\sinh\beta\theta+\left(\beta^2-\alpha^2\right)\cos\alpha\theta\cosh\beta\theta\right).
\end{align*}
Finally, recalling that $\theta(0)=0,$
\begin{align*}
\dddot u_1(0)&=\left(\ddot H(0)+H(0)^3\left(\beta^2-3\alpha^2\right)\right)\beta,\\ \dddot u_2(0)&=\left(\ddot H(0)+H(0)^3\left(3\beta^2-\alpha^2\right)\right)\alpha,\\  \dddot u_3(0)&=6H(0)\dot H(0)\alpha\beta.
\end{align*}
Therefore, 
\begin{equation*}
A=\begin{bmatrix}
	0 & 0 & 0 & 1 \\
	H(0)\beta & H(0)\alpha & 0 & \dot u_4(0) \\
	\dot H(0)\beta & \dot H(0)\alpha & 2H(0)^2\alpha\beta & \ddot u_4(0) \\
	\left(\ddot H(0)+H(0)^3\left(\beta^2-3\alpha^2\right)\right)\beta & \left(\ddot H(0)+H(0)^3\left(3\beta^2-\alpha^2\right)\right)\alpha & 6H(0)\dot H(0)\alpha\beta & \dddot u_4(0)
\end{bmatrix}.
\end{equation*}
In particular,
\begin{equation*}
\begin{split}
	\det A&=H(0)\beta\left(6H(0)\dot H(0)^2\alpha^2\beta-2H(0)^2\left(\ddot H(0)+H(0)^3\left(3\beta^2-\alpha^2\right)\right)\alpha^2\beta\right)\\
	&\quad-H(0)\alpha\left(6H(0)\dot H(0)^2\alpha\beta^2-2H(0)^2\left(\ddot H(0)+H(0)^3\left(\beta^2-3\alpha^2\right)\right)\alpha\beta^2\right)\\
	%    &\deleted{=H(0)^2\alpha^2\beta^2\left(6\dot H(0)^2-2H(0)\ddot H(0)-2H(0)^4\left(3\beta^2-\alpha^2\right)\right)}\\
	%    &\deleted{\quad-H(0)^2\alpha^2\beta^2\left(6\dot H(0)^2-2H(0)\ddot H(0)-2H(0)^4\left(\beta^2-3\alpha^2\right)\right)}\\
	&=-4H(0)^6\alpha^2\beta^2\left(\alpha^2+\beta^2\right).
\end{split}
\end{equation*}
Since $H(0),\alpha,\beta\neq 0$, we conclude that $\det A\neq 0$.
Assume instead that $\mu>\mu_0$ and $u_1,\ldots,u_4$ are the functions given by \eqref{soluzionicasoreale}. 
%that 
%  \begin{equation*}
%     C>-\frac{x_0}{L+x_0}+2\sqrt{\frac{x_0}{L+x_0}}.
% \end{equation*}
Set $z_1=\alpha-\gamma$ and $z_2=\alpha+\gamma$. %Let $x_1,z_2\in\mathbb C$ be as above. In this case, we recall that $z_1,z_2\in \rr$, with $z_1\neq z_2$, and that
Then $z_1,z_2\in\rr$, $z_1\neq z_2$ and
\begin{equation*}
u_1=\sin z_1 \theta,\qquad
u_2=\sin z_2\theta,\qquad
u_3= \cos z_1 \theta,\qquad
u_4=\cos z_2\theta.
\end{equation*}
Therefore
\begin{equation*}
\dot u_1=Hz_1\cos z_1 \theta,\qquad
\dot u_2=Hz_2\cos z_2\theta,\qquad
\dot u_3= -Hz_1\sin z_1 \theta,\qquad
\dot u_4=-Hz_2\sin z_2\theta.
\end{equation*}
Moreover,
\begin{align*}
\ddot u_1&=\dot Hz_1\cos z_1\theta-H^2z_1^2\sin z_1 \theta,\qquad \ddot u_2=\dot Hz_2\cos z_2\theta-H^2z_2^2\sin z_2 \theta,\\
\ddot u_3&=-\dot Hz_1\sin z_1\theta-H^2z_1^2\cos z_1 \theta,\qquad \ddot u_4=-\dot Hz_2\sin z_2\theta-H^2z_2^2\cos z_2 \theta.
\end{align*}
Finally,
\begin{align*}
\dddot u_1(0)&=\left(\ddot H(0)-H(0)^3z_1^2\right)z_1,\qquad  \dddot u_2(0)=\left(\ddot H(0)-H(0)^3z_2^2\right)z_2,\\
\dddot u_3(0)&=-3H(0)\dot H(0)z_1^2,\qquad  \dddot u_4(0)=-3H(0)\dot H(0)z_2^2.
\end{align*}
Therefore, 
\begin{equation*}
A=\begin{bmatrix}
	0 & 0 & 1 & 1 \\
	H(0)z_1 & H(0)z_2 & 0 & 0 \\
	\dot H(0)z_1 & \dot H(0)z_2 & -H(0)^2z_1^2 & -H(0)^2z_2^2 \\
	\left(\ddot H(0)-H(0)^3z_1^2\right)z_1 & \left(\ddot H(0)-H(0)^3z_2^2\right)z_2 & -3H(0)\dot H(0)z_1^2 & -3H(0)\dot H(0)z_2^2
\end{bmatrix}.
\end{equation*}
In particular,
\begin{equation*}
\begin{split}
	\det A &=H(0)z_1\left(-3H(0)\dot H(0)^2z_2^3+ H(0)^2\left(\ddot H(0)-H(0)^3z_2^2\right)z_2^3\right)\\
	&\quad- H(0)z_2\left(-3H(0)\dot H(0)^2z_1z_2^2+ H(0)^2\left(\ddot H(0)-H(0)^3z_1^2\right)z_1z_2^2\right)\\
	&\quad-H(0)z_1\left(-3H(0)\dot H(0)^2z_1^2z_2+ H(0)^2\left(\ddot H(0)-H(0)^3z_2^2\right)z_1^2z_2\right)\\
	&\quad+ H(0)z_2\left(-3H(0)\dot H(0)^2z_1^3+ H(0)^2\left(\ddot H(0)-H(0)^3z_1^2\right)z_1^3\right)\\
	%        &\deleted{=H(0)^2z_1z_2\left(-3\dot H(0)^2z_2^2+ H(0)\ddot H(0)z_2^2-H(0)^4z_2^4\right)}\\
	%      &\deleted{\quad- H(0)^2z_1z_2\left(-3\dot H(0)^2z_2^2+ H(0)\ddot H(0)z_2^2-H(0)^4z_1^2z_2^2\right)}\\
	%      &\deleted{\quad-H(0)^2z_1z_2\left(-3\dot H(0)^2z_1^2+ H(0)\ddot H(0)z_1^2-H(0)^4z_1^2z_2^2\right)}\\
	%      &\deleted{\quad+ H(0)^2z_1z_2\left(-3\dot H(0)^2z_1^2+ H(0)\ddot H(0)z_1^2-H(0)^4z_1^4\right)}\\
	&=-H(0)^6z_1z_2\left(z_1^2-z_2^2\right)^2.
\end{split}
\end{equation*}
Since $H(0),z_1,z_2\neq 0$ and $z_1\neq z_2$, we conclude that $\det A\neq 0$. Finally, assume that $\mu=\mu_0$ and $u_1,\ldots,u_4$ are the functions given by \eqref{soluzionicasocritico}. Recall that 
\begin{equation*}
u_1=\sin \alpha \theta,\qquad
u_2=\cos\alpha\theta,\qquad
u_3= \theta\sin\alpha \theta,\qquad
u_4=\theta\cos\alpha\theta.
\end{equation*}
Then,
\begin{equation*}
\dot u_1=\alpha H\cos\alpha\theta,\qquad\dot u_3=H\sin\alpha\theta+\alpha H\theta\cos\alpha\theta,\qquad \dot u_4 =H\cos\alpha\theta-\alpha H\theta\sin\alpha\theta.
\end{equation*}
Moreover,
\begin{align*}
\ddot u_1&=\alpha\dot H\cos\alpha\theta-\alpha^2H^2\sin\alpha\theta,\\
\ddot u_3&=\dot H\sin\alpha\theta+2\alpha H^2\cos\alpha\theta+\alpha\dot H\theta\cos\alpha\theta-\alpha^2 H^2\theta\sin\alpha\theta,\\
\ddot u_4&=\dot H\cos\alpha\theta-2\alpha H^2\sin\alpha\theta-\alpha\dot H\theta\sin\alpha\theta-\alpha^2H^2\theta\cos\alpha\theta.
\end{align*}
Finally,
\begin{align*}
\dddot u_1(0)&=\ddot H(0)\alpha-H(0)^3\alpha^3,\\
\dddot u_3(0)&=6H(0)\dot H(0)\alpha,\\
\dddot u_4(0)&=\ddot H(0)-3H(0)^3\alpha^2.
\end{align*}
Therefore, 
\begin{equation*}
A=\begin{bmatrix}
	0 & 1 & 0 & 0 \\
	H(0)\alpha & \dot u_2(0) & 0 & H(0) \\
	\dot H(0)\alpha & \ddot u_2(0) & 2H(0)^2\alpha & \dot H(0) \\
	\ddot H(0)\alpha-H(0)^3\alpha^3 & \dddot u_2(0) & 6H(0)\dot H(0)\alpha & \ddot H(0)-3H(0)^3\alpha^2
\end{bmatrix}.
\end{equation*}
In particular,
\begin{equation*}
\begin{split}
	\det A&=-H(0)\alpha\left(2H(0)^2\ddot H(0)\alpha-6H(0)^5\alpha^3-6H(0)\dot H(0)^2\alpha\right)\\
	&\quad-H(0)\left(6H(0)\dot H(0)^2\alpha^2-2H(0)^2\ddot H(0)\alpha^2+2H(0)^5\alpha^4\right)\\
	&=4 H(0)^6\alpha^4.
\end{split}
\end{equation*}
Since $H(0),\alpha\neq 0$, we conclude that $\det A\neq 0$.
\end{proof}

\end{document}
